\chapter{Ilustrações, Tabelas e Elementos Especiais}

Ilustrações, tabelas, gráficos, quadros, fluxogramas, equações e outros elementos especiais devem servir à compreensão do trabalho. Eles não devem ser inseridos apenas para preencher páginas ou tornar o texto visualmente mais chamativo.

\section{Ilustrações}

Ilustrações incluem figuras, gráficos, mapas, fotografias, desenhos, esquemas, fluxogramas e outros recursos visuais. Devem ser citadas no texto, posicionadas próximas ao trecho em que são discutidas e acompanhadas de identificação e fonte.

A identificação deve permitir que o leitor localize e entenda o item. A fonte informa origem ou autoria da imagem. Quando a ilustração for elaborada pelo próprio estudante, isso deve ser indicado.

A identificação aparece acima do elemento, com palavra designativa, número, travessão e título, por exemplo: ``Figura 3 -- Etapas da pesquisa''. Abaixo ficam a fonte e, quando necessárias, notas explicativas. O padrão institucional para autoria própria é: ``Fonte: elaboração própria.'' Essa fórmula deve ser usada de modo uniforme em figuras, quadros, gráficos, fluxogramas e demais ilustrações produzidas pela autoria do TCC.

\decisaoinstitucional

\section{Figuras comparativas}

Quando duas imagens são usadas para comparação, cada uma deve ter relação clara com o argumento. A simples colocação lado a lado não dispensa explicação no texto. O estudante deve indicar o que deve ser observado, quais diferenças importam e como elas sustentam a análise.

Se as imagens forem figuras independentes, cada uma recebe número, título, fonte e chamada próprios, mesmo que sejam posicionadas lado a lado. Use partes ``a'' e ``b'' sob uma única legenda apenas quando elas constituírem, de fato, uma figura composta inseparável.

\section{Tabelas e quadros}

Tabelas organizam dados, especialmente numéricos. Quadros costumam organizar informações textuais, categorias, comparações ou sínteses. Ambos devem ser legíveis, citados no texto e acompanhados de fonte.

As tabelas devem observar as Normas de apresentação tabular do IBGE. Seu número e título ficam acima; fonte e notas ficam abaixo. Quadros são tratados como ilustrações no padrão do trabalho e seguem a identificação institucional correspondente.

Evite tabelas muito largas, com excesso de linhas ou informação que poderia ser apresentada em texto. Quando necessário, dados extensos podem ser deslocados para apêndice.

\subsection{Diferença prática entre tabela e quadro}

Use tabela quando o foco estiver em dados numéricos, medidas, frequências, percentuais ou resultados organizados estatisticamente. Use quadro quando o foco estiver em texto, conceitos, categorias, critérios, sínteses ou comparação qualitativa.

Essa distinção ajuda a manter o padrão de apresentação e evita que todo elemento organizado em linhas e colunas seja chamado genericamente de tabela.

\subsection{Cuidados com tabelas}

Tabelas devem apresentar cabeçalho claro, unidades de medida, fonte e alinhamento que favoreça comparação. Não repita no texto todos os números da tabela; interprete os resultados mais importantes.

Quando uma tabela continuar em outra página, repita o cabeçalho e indique a continuidade com as expressões ``continua'', ``continuação'' e ``conclusão'', conforme a posição. Evite quebrar linhas de dados entre páginas. Unidades de medida devem estar no cabeçalho ou em posição que elimine ambiguidade.

\subsection{Cuidados com quadros}

Quadros devem ser usados para síntese textual. Evite quadros com parágrafos muito longos ou conteúdo que seria melhor explicado no corpo do texto.

\section{Gráficos}

Gráficos ajudam a visualizar tendências, proporções e comparações. Devem ter título claro, eixos identificados, unidades quando aplicáveis e fonte. Cores e estilos devem ser escolhidos para favorecer legibilidade, inclusive em impressão em preto e branco.

Não use apenas cor para distinguir séries. Combine cor com padrões, marcadores, rótulos ou estilos de linha e preserve contraste suficiente. Quando o arquivo for digital, forneça texto alternativo ou descrição equivalente no documento acessível.

\section{Fluxogramas}

Fluxogramas representam processos, etapas, decisões ou fluxos de trabalho. Devem usar símbolos de modo consistente e indicar início, fim, atividades e decisões quando representarem um processo completo.

No padrão institucional, fluxogramas podem ser tratados como ilustrações quando aparecem em pequena quantidade. Quando forem recorrentes no trabalho, pode ser adotada lista própria para facilitar consulta.

A ABNT NBR 14724:2024 disciplina a apresentação do fluxograma como ilustração, mas não fixa uma gramática completa de símbolos. O trabalho deve adotar convenção reconhecida na área, explicá-la quando necessário e manter consistência. Símbolos, setas e decisões precisam ser legíveis e não podem depender exclusivamente de cores.

\section{Equações}

Equações devem ser numeradas quando forem citadas posteriormente ou quando houver necessidade de referência. Variáveis e símbolos precisam ser definidos no texto ou em lista própria, se o trabalho utilizar muitos símbolos.

\section{Pseudocódigo}

Pseudocódigo pode ser usado em trabalhos que descrevem algoritmos, rotinas, processos computacionais ou procedimentos lógicos. Deve ser legível para público acadêmico e não depender de detalhes de uma linguagem de programação específica, salvo quando esse for o objeto do trabalho.

Apresente o pseudocódigo como algoritmo ou elemento técnico numerado, com título, fonte, entradas, saídas e passos claramente identificados. Numeração de linhas e palavras-chave em destaque ajudam a discussão. Se o bloco não couber na página, prefira movê-lo integralmente, reduzir espaços ou dividi-lo em partes logicamente nomeadas; não reduza a fonte a ponto de prejudicar a leitura.

\caixaorientacao{Todo elemento visual deve ser comentado no texto. A frase ``conforme figura a seguir'' é insuficiente se o estudante não explicar o que a figura demonstra.}
